\documentclass[11pt,a4paper]{article}
\usepackage[utf8]{inputenc}
\usepackage[T1]{fontenc}
\usepackage{lmodern}
\usepackage[a4paper,margin=0.75in]{geometry}
\usepackage{microtype}
\usepackage{amsmath,amssymb,amsthm}
\usepackage{booktabs}
\usepackage{tabularx}
\usepackage[numbers]{natbib}
\usepackage{algorithm}
\usepackage{algorithmic}

\newtheorem{theorem}{Theorem}
\newtheorem{proposition}{Proposition}
\newtheorem{corollary}{Corollary}
\newtheorem{remark}{Remark}
\newtheorem{definition}{Definition}
\newcolumntype{Y}{>{\raggedright\arraybackslash}X}

\title{Graph Energy Transformer with Explicit Motif Retrieval}
\author{}
\date{}

\setlength{\parindent}{0pt}
\setlength{\parskip}{0.5em}

\begin{document}
\maketitle

\begin{abstract}
We introduce the \emph{Graph Energy Transformer} (GET), a graph model whose
scalar objective combines graph-local pairwise retrieval, sparse order-3 motif
retrieval over anchored wedges and triangles, ET-style global attention, and a
Hopfield/FFN memory term. GET is intended as an energy formulation of the main
interaction modules of a modern graph transformer: local message passing,
global mixing, higher-order motif retrieval, and tokenwise feed-forward memory.
It reduces to continuous Modern Hopfield retrieval in a single-node memory-only
limit and to a restricted ET attention-plus-HN energy in the complete-support
limit. The log-sum-exp retrieval terms admit a standard entropy-regularized
simplex lifting, so the softmax weights can be interpreted as auxiliary routing
variables. This identity is not itself novel; its role is to make the update
field precise. Differentiating one scalar energy gives both anchor-side
aggregation and reverse-incidence terms, as in ET, rather than an arbitrary
feed-forward message-passing rule. The motif branch uses a rank-$R$
factorization that is not representable by the additive-separable pairwise
score used by the local branch. Under fixed local, motif, and global supports
and $\varepsilon$-stabilized LayerNorm, the energy is smooth and coercive,
accepted Armijo steps decrease it monotonically, and strict local minima admit
a local Lyapunov certificate. The empirical claim is deliberately narrow:
FullGET should improve over a capacity-matched no-motif GET harness on tasks
whose signal is local and motif-causal.
\end{abstract}

%----------------------------------------------------------------------
\section{Introduction}
%----------------------------------------------------------------------

Energy-based models define inference by differentiating a scalar objective.
That constraint is stronger than ordinary message passing: the update field
must be an integrable gradient field, and a node receives terms both from
queries it anchors and from other queries in which it appears as a key.
Energy Transformer (ET) makes this distinction explicit for self-attention
\citep{hoover2023energy}; Modern Hopfield Networks provide the associated
log-sum-exp retrieval form \citep{ramsauer2021hopfield}.

GET asks a narrower graph question. If a modern graph-transformer harness is
written as a Hopfield/ET-style energy over graph states, can an explicit sparse
order-3 motif term be isolated and tested against the same energy harness with
the motif branch removed? This is not a claim that higher-order graph
reasoning, graph attention, local/global graph transformers, or energy-based
graph dynamics are new. Those are already covered by
MPNNs and GATs \citep{gilmer2017neural,velickovic2018graph}, graph
transformers \citep{ying2021graphormer,rampasek2022gps,ma2023grit,
shirzad2023exphormer}, gradient-flow and diffusion views of GNNs
\citep{digiovanni2023gnnflow,wu2025diffusion}, Graph Hopfield Networks
\citep{rao2026graph}, and higher-order/substructure models
\citep{morris2019weisfeiler,bouritsas2023subgraph,bodnar2021mpsn,
hussain2024triplet}.

The purpose of GET is therefore methodological and empirical: define one scalar
graph energy with local, motif, global, and Hopfield/FFN retrieval terms;
derive the full gradient update rather than postulate a feed-forward
aggregator; and test whether the motif term matters under matched no-motif
controls.

\paragraph{Inheritance chain.}
GET is built so that known ancestors appear as special cases by disabling
inactive branches while keeping active couplings nonnegative
(Table~\ref{tab:special-cases}). The two anchor limits are a Modern Hopfield
memory-only case and a restricted complete-support ET-style attention-plus-HN
case; the full model adds graph-local and sparse motif retrieval on top of the
global ET-style harness.

\begin{table}[htbp]
\centering
\caption{Special cases of GET. $\lambda_2$, $\lambda_3$, $\lambda_g$, and
$\lambda_m$ are the local pairwise, motif, global attention, and HN/FFN memory
couplings.}
\label{tab:special-cases}
\renewcommand{\arraystretch}{1.2}
\begin{tabular}{lcccc l}
\toprule
Setting & $\lambda_2$ & $\lambda_3$ & $\lambda_g$ & $\lambda_m$ & Reduces to \\
\midrule
$|V|=1$, no edges & $0$ & $0$ & $0$ & $>0$ & Modern Hopfield Network \citep{ramsauer2021hopfield} \\
Complete global support & $0$ & $0$ & $>0$ & $>0$ & restricted ET attention+HN \citep{hoover2023energy} \\
Sparse graph only & $>0$ & $0$ & $0$ & $0$ & graph-local pairwise retrieval \\
No-motif harness & $>0$ & $0$ & $>0$ & $>0$ & local/global/HN graph transformer \\
Full model & $>0$ & $>0$ & $>0$ & $>0$ & GET \\
\bottomrule
\end{tabular}
\end{table}

\paragraph{Contributions.}
\begin{enumerate}
\item We define a scalar graph energy with local pairwise retrieval, explicit
  sparse motif retrieval, ET-style global attention, and a Hopfield/FFN memory
  branch. The motif score is an explicit sparse order-3 interaction that is
  generically not representable by the additive-separable pairwise family used
  by the graph-local branch.
\item We verify the inheritance chain to Modern Hopfield retrieval and to a
  restricted ET-style attention-plus-HN energy. We also make explicit the
  entropy-regularized simplex lifting of the log-sum-exp terms; this identifies
  the softmax weights as routing variables, not as Bethe or VMP beliefs.
\item We show that full-gradient descent yields region-message dynamics with
  both anchor-side and reverse-incidence terms over local edges, retained
  motifs, global attention supports, and tokenwise memory. This is the main
  structural difference from a plain feed-forward MPNN, GAT, or graph
  transformer update.
\item Under fixed local, motif, and global supports and
  $\varepsilon$-stabilized LayerNorm, we show coercivity, local smoothness,
  monotone decrease under accepted Armijo steps, and local Lyapunov stability
  around strict local minima.
\item We organize evaluation around mechanism identification first and transfer
  second: motif-causal diagnostics, realized-expressivity tests such as BREC,
  bounded-degree molecular graphs, and only then long-range or ET-style graph
  benchmarks.
\end{enumerate}

%----------------------------------------------------------------------
\section{Background and Special Cases}
%----------------------------------------------------------------------

\subsection{Modern Hopfield Networks}

Given a query $\xi\in\mathbb{R}^d$ and stored patterns
$\{b_a\}_{a=1}^K\subset\mathbb{R}^d$, the Modern Hopfield energy
\citep{ramsauer2021hopfield} is
\begin{equation}
  E_{\mathrm{Hop}}(\xi)
  = \frac{1}{2}\|\xi\|^2
    - \frac{1}{\beta}\log\sum_{a=1}^K\exp\!\bigl(\beta\,\xi^\top b_a\bigr)
    + C(\beta,\{b_a\}),
  \label{eq:mhn-scalar}
\end{equation}
where $C$ is independent of $\xi$. Stationarity gives the retrieval fixed
point
\begin{equation}
  \xi^* = B\,\mathrm{softmax}(\beta B^\top \xi^*),
  \qquad B=[b_1,\ldots,b_K].
  \label{eq:mhn-fixedpt}
\end{equation}

\subsection{Energy Transformer}

For token states $X\in\mathbb{R}^{N\times d}$, the attention part of the ET
energy \citep{hoover2023energy} is
\begin{equation}
  E_{\mathrm{ET}}^{\mathrm{ATT}}(X)
  = \frac{1}{2}\|X\|_F^2
    - \frac{1}{\beta}\sum_{i=1}^N
      \log\sum_{j=1}^N\exp\!\bigl(\beta s_{ij}\bigr),
  \qquad
  s_{ij} = \frac{1}{\sqrt{d}}\,{(W_Q x_i)}^\top(W_K x_j).
  \label{eq:et}
\end{equation}
GET recovers this restricted single-head score model in the complete-support
global-attention limit with matched support, scaling, and normalization
conventions.

The published ET block also contains a Hopfield Network module
$E_{\mathrm{ET}}^{\mathrm{HN}}$, applied in parallel with attention, whose
weights act as memories and whose tied projection structure is the ET analogue
of a transformer feed-forward network. GET adopts the same decomposition at
the graph level: local graph attention, motif attention, global attention, and
HN/FFN memory are separate additive terms in one scalar energy. This is a
single-state ET-style graph energy, not a full hierarchical associative memory
over all layers; the HAM viewpoint is useful lineage, but the guarantee here is
descent of the implemented scalar energy \citep{hoover2022universal}.

The important ET lesson for this paper is not only the scalar energy, but the
gradient. Let
$a_{ij}=\mathrm{softmax}_{j}(\beta s_{ij})$. Ignoring normalization constants
and writing only the attention contribution, the negative gradient contains
\begin{equation}
  -\nabla_{x_i}E_{\mathrm{ET}}
  \supset
  W_Q^\top\sum_j a_{ij}W_Kx_j
  +
  W_K^\top\sum_u a_{ui}W_Qx_u .
  \label{eq:et-two-term}
\end{equation}
The first term is query-side aggregation. The second appears because $x_i$ is
also a key in other tokens' scores. A feed-forward attention layer may keep only
the first effect; exact energy descent does not. The same two-term structure
appears for GET's local and global attention branches.

\subsection{Message Passing and Free Energy}

Classical BP, GBP, and VMP already show that message passing can be tied to
free-energy stationarity or variational coordinate optimization
\citep{yedidia2005free,winn2005vmp}. This does not make GET a probabilistic
inference algorithm. BP and GBP optimize constrained beliefs with marginal
consistency; VMP optimizes a factorized posterior bound. GET instead optimizes
continuous node states, while its auxiliary weights are routing distributions
over graph supports. The overlap with graphical-model inference is algebraic,
not semantic. Surveys of the GNN--PGM intersection make the same distinction
between learned representation passing and probabilistic belief passing
\citep{hua2022gnnpgm}.

The only safe region-graph analogy is structural: local pairwise scores act on
edge regions, motif scores act on retained triangle or wedge regions, global
attention scores act on long-range node-pair regions, and memory scores act on
node-to-prototype regions. GET does not introduce region beliefs, Möbius
counting numbers, or marginal-consistency constraints, so it should not be
identified with Kikuchi free energy or GBP.

\subsection{Positioning}

The narrow contribution is an ET-style graph energy harness with an explicit
sparse motif term. Higher-order GNNs, substructure counting,
simplicial message passing, relation-aware transformers, and Triplet Graph
Transformer already cover motifs, tuples, and higher-order graph interactions
\citep{morris2019weisfeiler,bouritsas2023subgraph,bodnar2021mpsn,
diao2023relational,hussain2024triplet,burns2023simplicial}. GraphGPS,
Graphormer, GRIT, Exphormer, GPSE, and related graph transformers cover
structural encodings, local/global attention, and scalable transformer
harnesses \citep{ying2021graphormer,rampasek2022gps,ma2023grit,
shirzad2023exphormer,canturk2024gpse,airale2025spse,bao2025homomorphism}.
Graph Hopfield Networks are especially close in lineage, but their graph term is
closer to Hopfield retrieval coupled with graph smoothing for node
classification; GET's distinctive term is learned motif-indexed retrieval inside
the same local/global/HN energy harness.
GET should be judged on the smaller question: does a motif term inside one
scalar graph-transformer energy improve over the matched no-motif harness?

%----------------------------------------------------------------------
\section{Graph Energy Transformer}
%----------------------------------------------------------------------

\subsection{State Space and Support}

Let $G=(V,E)$ be an undirected graph with hidden dimension $d\ge 2$. Each
node $i\in V$ has raw state $x_i\in\mathbb{R}^d$ and normalized state
\begin{equation}
  g_i = \mathrm{LN}(x_i),
  \label{eq:ln}
\end{equation}
where LayerNorm is stabilized with $\varepsilon>0$. Let $\mathcal{N}(i)$ be
the pairwise support of node $i$ (optionally including a self-loop), and let
$\mathcal{W}(i)$ be a retained set of anchored unordered neighbor pairs around
$i$. We write $(j,k)\in\mathcal W(i)$ for a representative of the unordered pair
and store each pair once. The intended motifs are sparse open wedges and
triangles, though ordered or typed supports can be used when role-asymmetric
patterns are required. Let $\mathcal S_g(i)$ be the global attention support
of node $i$, such as all nodes, a CLS-mediated support, anchors, or a sparse
long-range support. Finally, let $\mathcal{B}=\{b_a\}_{a=1}^{K}$ denote the
HN/FFN memory bank.

For theory, the key assumption is that all supports are fixed during one
inference trajectory: $\mathcal{W}(i)$ and $\mathcal S_g(i)$ are chosen once by
an invariant rule such as exact wedge/triangle support, a deterministic score
threshold, a bounded motif budget, dense global support, or a fixed sparse
long-range pattern. That assumption keeps the energy smooth between support
changes and makes the descent results interpretable as properties of a fixed
graph retrieval problem rather than of a changing support selector.

\subsection{Local, Motif, Global, and HN/FFN Scores}

The pairwise branch uses linear query--key maps $q_2,k_2$ and score
\begin{equation}
  s^{(2)}_{ij}
  = \frac{1}{\sqrt d}\,q_2(g_i)^\top k_2(g_j) + a^{(2)}_{ij},
  \qquad j\in\mathcal{N}(i),
  \label{eq:s2}
\end{equation}
where $a^{(2)}_{ij}$ is an optional edge-feature bias.
Unless $q_2=k_2$ and the biases are symmetric, $s^{(2)}_{ij}$ need not equal
$s^{(2)}_{ji}$. Thus GET is not an undirected graphical-model energy on edges;
it is an attention-style scalar energy over anchored retrieval problems.

The motif branch uses a rank-$R$ factorization with motif type embedding
$t_{\tau_{ijk}}^{(r)}\in\mathbb{R}^{d}$:
\begin{equation}
  s^{(3)}_{i,jk}
  = \frac{1}{\sqrt{Rd}}\sum_{r=1}^{R}
      {q_3^{(r)}(g_i)}^\top
      \bigl(k_3^{(r)}(g_j)\odot k_3^{(r)}(g_k) + t_{\tau_{ijk}}^{(r)}\bigr),
  \qquad (j,k)\in\mathcal{W}(i).
  \label{eq:s3}
\end{equation}
Because the key term is multiplicative in $(g_j,g_k)$, this branch captures
joint neighbor effects that additive pairwise scores cannot express in
general. A triangle may contribute once for each anchor whose support retains
the other two vertices; the motif energy is therefore anchor-local, not a single
shared Kikuchi-style triangle potential.

The global branch uses its own query--key maps $q_g,k_g$:
\begin{equation}
  s^{(g)}_{ij}
  = \frac{1}{\sqrt d}\,q_g(g_i)^\top k_g(g_j) + a^{(g)}_{ij},
  \qquad j\in\mathcal S_g(i).
  \label{eq:sg}
\end{equation}
This is the ET-style graph-transformer mixing term. In small graphs
$\mathcal S_g(i)=V$ gives dense attention; in larger graphs it can be a fixed
anchor, expander, sampled, or CLS-mediated support. The scalar-energy gradient
again includes both query-side and reverse key-side contributions.

The HN/FFN branch follows the Modern Hopfield / ET pattern with linear maps
$q_m$ and $k_m$:
\begin{equation}
  s^{(m)}_{i,a}
  = \frac{1}{\sqrt d}\,q_m(g_i)^\top k_m(b_a),
  \qquad a\in\{1,\dots,K\}.
  \label{eq:sm}
\end{equation}

\subsection{Why the Motif Term Is Genuinely Higher-Order}

The motif score in~\eqref{eq:s3} is not merely a rephrased pairwise model.
Suppose one tries to represent the dependence on $(g_j,g_k)$ by an
additive-separable family of the form
$u(g_i,g_j)+v(g_i,g_k)+c_{\tau_{ijk}}$. Such a score has zero mixed second
derivative with respect to any coordinate of $g_j$ and any coordinate of
$g_k$. By contrast, the term
${q_3^{(r)}(g_i)}^\top\bigl(k_3^{(r)}(g_j)\odot k_3^{(r)}(g_k)\bigr)$ has
generically nonzero cross-dependence because the two neighbor states enter
multiplicatively. Except in degenerate parameter settings, it therefore
cannot be rewritten as a sum of independent edgewise contributions.

This is the sense in which GET adds an explicit order-3 interaction rather
than a richer pairwise scorer. The contribution of neighbor $j$ to anchor $i$
depends on which other neighbor $k$ is paired with it, so the model can
distinguish local configurations that match on degree and pairwise marginals
but differ in wedge closure or triangle structure.
This nonseparability claim is only relative to the additive edgewise score
family in the no-motif local branch. It is not a theorem against arbitrary
higher-order or pair-state architectures.

\subsection{Block Energy and Inference}

The attention energy combines the quadratic penalty with graph-local pairwise,
motif, and global retrieval:
\begin{equation}
\begin{aligned}
  E_{\mathrm{att}}(X)
  ={}& \frac{1}{2}\sum_{i\in V}\|x_i\|^2
     - \frac{\lambda_2}{\beta_2}\sum_{i\in V}
       \log\sum_{j\in\mathcal{N}(i)} \exp\!\bigl(\beta_2 s^{(2)}_{ij}\bigr) \\
    &- \frac{\lambda_3}{\beta_3}\sum_{i\in V}
       \log\sum_{(j,k)\in\mathcal{W}(i)}
         \exp\!\bigl(\beta_3 s^{(3)}_{i,jk}\bigr)\\
    &- \frac{\lambda_g}{\beta_g}\sum_{i\in V}
       \log\sum_{j\in\mathcal S_g(i)}
         \exp\!\bigl(\beta_g s^{(g)}_{ij}\bigr),
\end{aligned}
  \label{eq:eatt}
\end{equation}
with the convention that empty supports contribute zero. The HN/FFN memory
branch is
\begin{equation}
  E_{\mathrm{HN}}(X)
  = -\frac{\lambda_m}{\beta_m}\sum_{i\in V}
      \log\sum_{a=1}^{K}\exp\!\bigl(\beta_m s^{(m)}_{i,a}\bigr),
  \label{eq:ehn}
\end{equation}
and the full block energy is
\begin{equation}
  E_{\mathrm{GET}}(X) = E_{\mathrm{att}}(X) + E_{\mathrm{HN}}(X).
  \label{eq:fullenergy}
\end{equation}

\subsection{Energy-Induced Message Weights}

The log-sum-exp form has a standard variational identity. It is useful here
only because it fixes what the local ``messages'' are; by itself, this identity
is not a novelty claim.

\begin{proposition}[Entropy-regularized routing form]
\label{prop:routing}
For fixed $X$ and positive temperatures, $E_{\mathrm{GET}}(X)$ equals
$\min_{\alpha,\gamma,\pi,\rho}\widetilde E(X,\alpha,\gamma,\pi,\rho)$ over
local simplexes
$\alpha_i\in\Delta_{\mathcal N(i)}$,
$\gamma_i\in\Delta_{\mathcal W(i)}$,
$\pi_i\in\Delta_{\mathcal S_g(i)}$, and
$\rho_i\in\Delta_K$, where
\begin{equation}
\begin{aligned}
\widetilde E
={}& \frac12\sum_i\|x_i\|^2
-\lambda_2\sum_i\sum_{j\in\mathcal N(i)}\alpha_{ij}s^{(2)}_{ij}
-\frac{\lambda_2}{\beta_2}\sum_i H(\alpha_i)\\
&-\lambda_3\sum_i\sum_{(j,k)\in\mathcal W(i)}
  \gamma_{i,jk}s^{(3)}_{i,jk}
-\frac{\lambda_3}{\beta_3}\sum_i H(\gamma_i)\\
&-\lambda_g\sum_i\sum_{j\in\mathcal S_g(i)}\pi_{ij}s^{(g)}_{ij}
-\frac{\lambda_g}{\beta_g}\sum_i H(\pi_i)\\
&-\lambda_m\sum_i\sum_a\rho_{ia}s^{(m)}_{ia}
-\frac{\lambda_m}{\beta_m}\sum_i H(\rho_i).
\end{aligned}
\label{eq:lifted-energy}
\end{equation}
The minimizers are the softmax weights
\begin{equation}
\begin{aligned}
  \alpha_{ij}=\mathrm{softmax}_{j\in\mathcal N(i)}(\beta_2s^{(2)}_{ij}),
  \qquad
  \gamma_{i,jk}=\mathrm{softmax}_{(j,k)\in\mathcal W(i)}
    (\beta_3s^{(3)}_{i,jk}),\\
  \pi_{ij}=\mathrm{softmax}_{j\in\mathcal S_g(i)}(\beta_gs^{(g)}_{ij}),
  \qquad
  \rho_{ia}=\mathrm{softmax}_{a}(\beta_ms^{(m)}_{ia}).
\end{aligned}
\end{equation}
\end{proposition}

\begin{proof}
Use
$(1/\beta)\log\sum_r e^{\beta s_r}
=\max_{p\in\Delta}\{\sum_r p_rs_r+\beta^{-1}H(p)\}$ independently for the
edge, motif, global, and memory supports, then multiply each term by its
negative coupling.
\end{proof}

These variables are routing weights, not calibrated beliefs. The pairwise and
motif branches can be viewed as learned neural factors on edge and motif
regions, but GET inherits only the convex-analytic shape of free-energy message
passing. It does not inherit BP, GBP, VMP, or Kikuchi semantics unless one adds
belief variables, consistency constraints, region counting numbers, and an
explicit probabilistic model.

Inference minimizes $E_{\mathrm{GET}}$ over node states with parameters fixed:
\begin{equation}
  X^{t+1} = X^t - \eta_t\nabla_X E_{\mathrm{GET}}(X^t).
  \label{eq:update}
\end{equation}
For monotone inference we choose $\eta_t$ by Armijo backtracking. Starting
from $\eta_0$, shrink by $\gamma\in(0,1)$ until
\begin{equation}
  E_{\mathrm{GET}}(X^t - \eta_t G^t)
  \le E_{\mathrm{GET}}(X^t) - c\eta_t\|G^t\|_F^2,
  \qquad G^t = \nabla_X E_{\mathrm{GET}}(X^t),
  \label{eq:armijo}
\end{equation}
with $c\in(0,1)$.

\begin{minipage}{\textwidth}
\begin{algorithm}[H]
\caption{GET Inference by Monotone Armijo Descent}
\label{alg:armijo}
\begin{algorithmic}[1]
\STATE \textbf{Input:} initial state $X^0$, energy $E(\cdot)$, $c,\gamma\in(0,1)$, initial step $\eta_0$, steps $T$
\FOR{$t=0$ \textbf{to} $T-1$}
  \STATE $G^t \gets \nabla_X E(X^t)$
  \STATE $\eta \gets \eta_0$
  \WHILE{$E(X^t-\eta G^t) > E(X^t) - c\eta\|G^t\|_F^2$}
    \STATE $\eta \gets \gamma\eta$
  \ENDWHILE
  \STATE $X^{t+1} \gets X^t - \eta G^t$
\ENDFOR
\STATE \textbf{Return:} $X^T$
\end{algorithmic}
\end{algorithm}
\end{minipage}

The gradient derivation is given in Appendix~\ref{app:gradients}. As in
ET~\eqref{eq:et-two-term}, the local and global attention branches give both
anchor-side aggregation and reverse-incidence terms from other anchors for
which the node is a key. The motif branch similarly gives anchor motif messages
and reverse motif-incidence messages. This full-gradient field is the object
being studied; an anchor-only attention approximation would no longer be exact
descent on~\eqref{eq:fullenergy}. To keep the mechanism identifiable, score
maps remain linear and nonlinear capacity is confined to the input encoder and
task head.

\subsection{Input Encoding, Readout, and Training}

To keep the mechanism identifiable, the energy maps
$q_2,k_2,q_3,k_3,q_g,k_g,q_m,k_m$ remain linear and nonlinear capacity is
confined to the input encoder and task heads. Raw node attributes may be embedded or
projected by a small MLP, with optional degree, Laplacian, or random-walk
encodings when those same structural features are shared across all methods in
a comparison. Edge attributes enter the pairwise branch through the scalar
bias $a^{(2)}_{ij}$, while typed motifs use the embedding $t_{\tau_{ijk}}$.
The global branch uses either dense support or a fixed sparse long-range
support. The couplings $\lambda_2$, $\lambda_3$, $\lambda_g$, and $\lambda_m$
act as nonnegative branch gates and should be learned without weight decay.
For implementation, local pairwise and motif energies should be evaluated in
sparse edge-index form with grouped segment/scatter log-sum-exp reductions,
not with dense $N\times N$ attention masks. The global branch can use dense
attention on small graphs and sparse/anchor support on large graphs.

After $T$ inference steps, node-level tasks read out
$z_i=\mathrm{LN}(x_i^{(T)})$ directly. Graph-level tasks may use a simple
permutation-invariant pooled representation such as
\begin{equation}
  r_G = \Bigl[\tfrac{1}{|V|}\sum_i z_i\ ;\ \sum_i z_i\ ;\ \max_i z_i\Bigr],
\end{equation}
followed by an MLP head. When a benchmark genuinely rewards ET-style graph
wrappers, a CLS state can be included as an additional token in the same global
attention and HN/FFN energy and serve as the graph readout. Training should
distinguish the inner inference loop from the outer parameter
update:~\eqref{eq:update} refines latent states for fixed parameters, while
AdamW or another stochastic optimizer updates the parameters of the encoder,
score maps, memory bank, and task head through the unrolled solver. Fixed-step
training is practical; Armijo backtracking is most important for
evaluation-time inference and implementation checks.

%----------------------------------------------------------------------
\section{Theoretical Properties}
%----------------------------------------------------------------------

\subsection{Special-Case Reductions}

\begin{proposition}[Single-node case: Modern Hopfield retrieval]
\label{prop:single-node}
Let $|V|=1$ and set $\lambda_2=\lambda_3=\lambda_g=0$, $K\ge 1$. Then
\begin{equation}
  E_{\mathrm{GET}}(x)
  = \tfrac{1}{2}\|x\|^2
    - \tfrac{\lambda_m}{\beta_m}\log\sum_{a=1}^{K}
      \exp\!\bigl(\beta_m s^{(m)}_{1,a}\bigr).
\end{equation}
Setting $\nabla_x E_{\mathrm{GET}}=0$ gives
\begin{equation}
  x^* = \frac{\lambda_m}{\sqrt d}\,{J_{\mathrm{LN}}(x^*)}^\top
        {J_{q_m}(g^*)}^\top\sum_{a=1}^{K}\rho_a(x^*)\,k_m(b_a),
  \label{eq:single-node-fixedpt}
\end{equation}
where $g^*=\mathrm{LN}(x^*)$ and
$\rho_a(x)=\mathrm{softmax}_a(\beta_m s^{(m)}_{1,\cdot})$. If
$q_m=k_m=\mathrm{Id}$ and $\mathrm{LN}=\mathrm{Id}$, then choosing
$\lambda_m=\sqrt d$ and $\beta_m=\beta\sqrt d$ recovers the standard Modern
Hopfield retrieval rule
$x^*=\sum_a \mathrm{softmax}_a(\beta x^{*\top}b_a)\,b_a$.
\end{proposition}

\begin{proof}
Under the stated restrictions the energy contains only the quadratic term and
the memory log-sum-exp term. Differentiating with respect to $x$ yields the
stationarity equation above; the identity-map, no-LayerNorm limit then matches
\eqref{eq:mhn-fixedpt} after the stated scaling choice.
\end{proof}

\begin{proposition}[Complete-support case: restricted ET-style attention+HN]
\label{prop:et}
Let $\mathcal S_g(i)=V$ for all $i$, and set
$\lambda_2=\lambda_3=0$, $\lambda_g=1$, $q_g=W_Q$, and $k_g=W_K$. Then the
global attention branch of $E_{\mathrm{GET}}$ coincides with the single-head
complete-support ET attention energy in~\eqref{eq:et} up to the LayerNorm
reparameterization used in this paper. Keeping $\lambda_m>0$ adds the softmax
Modern Hopfield version of ET's HN/FFN memory term. This is still a restricted
ET reduction: it omits multi-head details, ET's more general activation
choices in $E_{\mathrm{ET}}^{\mathrm{HN}}$, and implementation-specific
wrappers.
\end{proposition}

\begin{proof}
With complete global support, $\mathcal S_g(i)=V$ for every node, so the
global log-sum-exp ranges over all token pairs. With local and motif branches
disabled, the global branch reduces exactly to the ET score model when
$\mathrm{LN}=\mathrm{Id}$ and otherwise differs only by the normalized-state
parameterization. The memory branch is the log-sum-exp Modern Hopfield
specialization of the ET HN module.
\end{proof}

\begin{proposition}[Fixed-point self-consistency]
\label{prop:fixedpt}
A stationary point $X^*$ of $E_{\mathrm{GET}}$ satisfies, for every $i\in V$,
\begin{equation}
  x_i^* = -{J_{\mathrm{LN}}(x_i^*)}^\top
           \nabla_{g_i} E_{\mathrm{GET}}(X^*).
  \label{eq:fixedpt-full}
\end{equation}
Thus every fixed point is a self-consistent retrieval field over graph-local
pairwise context, retained motifs, global attention context, and HN/FFN memory.
\end{proposition}

Because a node may also appear as a key in other nodes' local, motif, or global
scores, differentiating the scalar energy produces reverse-incidence terms: a
node receives gradient contributions from other nodes that attend to it. That
reciprocal flow is a structural consequence of the scalar-energy formulation
and is part of what distinguishes GET from a purely feed-forward attention
block.

\subsection{Descent and Stability}

\begin{proposition}[Coercivity and existence of minimizers]
\label{prop:coercive}
Assume fixed local, motif, and global supports; finite $K$; finite support
sizes; linear score maps; $d\ge 2$; and $\varepsilon$-stabilized LayerNorm. Then
$E_{\mathrm{GET}}$ is coercive and attains a global minimizer.
\end{proposition}

\begin{proof}[Proof sketch]
With $\varepsilon>0$, LayerNorm maps each raw state to a uniformly bounded
normalized state. Without affine parameters, $\|\mathrm{LN}(x_i)\|\le\sqrt d$;
with fixed affine scale $a$ and shift $b$,
$\|\mathrm{LN}(x_i)\|\le \|a\|_\infty\sqrt d+\|b\|$. On fixed support, the
local, motif, global, and memory scores are therefore uniformly bounded for
fixed parameters. The negative log-sum-exp branches are bounded below by a
constant, while the quadratic penalty
$\frac{1}{2}\sum_i\|x_i\|^2$ diverges as $\|X\|_F\to\infty$. Hence
$E_{\mathrm{GET}}(X)\to\infty$ as $\|X\|_F\to\infty$, and continuity yields
existence of a minimizer on a large sublevel set.
\end{proof}

\begin{proposition}[Local Lipschitz continuity of the gradient]
\label{prop:lipschitz}
Assume $\varepsilon>0$ in LayerNorm and linear learnable maps. Then
$\nabla_X E_{\mathrm{GET}}$ is locally Lipschitz on
$\mathbb{R}^{|V|\times d}$.
\end{proposition}

\begin{proof}[Proof sketch]
$E_{\mathrm{GET}}$ is built from smooth operations: linear maps,
$\varepsilon$-stabilized LayerNorm, Hadamard products, dot products, and
$\log\!\sum\exp$. Finite compositions of smooth maps are smooth, so the energy
is $C^1$ and its gradient is locally Lipschitz.
\end{proof}

\begin{theorem}[Armijo steps exist; energy is monotone]
\label{thm:monotone}
Under the conditions of Proposition~\ref{prop:lipschitz}, for every iterate
$X^t$ and every choice of $c\in(0,1)$, $\gamma\in(0,1)$, and $\eta_0>0$, the
backtracking loop terminates at some $\eta_t>0$ satisfying~\eqref{eq:armijo}.
Consequently,
\begin{equation}
  E_{\mathrm{GET}}(X^{t+1})
  \le E_{\mathrm{GET}}(X^t) - c\eta_t\|\nabla_X E_{\mathrm{GET}}(X^t)\|_F^2
  \le E_{\mathrm{GET}}(X^t).
\end{equation}
\end{theorem}

\begin{proof}
By Proposition~\ref{prop:lipschitz}, the gradient is locally Lipschitz near
$X^t$, so the standard descent lemma applies. For sufficiently small $\eta$,
one has $E(X^t-\eta G^t)\le E(X^t)-c\eta\|G^t\|_F^2$. Backtracking reaches
that regime after finitely many reductions by $\gamma$, which proves
existence of an accepted step and monotone decrease.
\end{proof}

\begin{proposition}[Local Lyapunov stability of strict local minima]
\label{prop:lyapunov}
Assume the conditions of Propositions~\ref{prop:coercive}
and~\ref{prop:lipschitz}, and let $X^*$ be a strict local minimizer of
$E_{\mathrm{GET}}$. Define the shifted energy
\begin{equation}
  V(X)=E_{\mathrm{GET}}(X)-E_{\mathrm{GET}}(X^*).
  \label{eq:lyapunov}
\end{equation}
Then $V$ is locally positive definite around $X^*$ and nonincreasing along
both the gradient flow $\dot X=-\nabla_X E_{\mathrm{GET}}(X)$ and the accepted
Armijo iteration. Hence $V$ is a local Lyapunov function for the inference
dynamics, and $X^*$ is Lyapunov stable.
\end{proposition}

\begin{proof}[Proof sketch]
Strict local minimality makes $V$ locally positive definite. Along gradient
flow, $\dot V(X)=-\|\nabla_X E_{\mathrm{GET}}(X)\|_F^2\le 0$ by the chain
rule. Along accepted Armijo steps, Theorem~\ref{thm:monotone} yields the
analogous discrete decay inequality. The full invariant-neighborhood argument
is recorded in Appendix~\ref{app:lyapunov}.
\end{proof}

\begin{remark}
These guarantees concern deterministic inference with fixed parameters. They
do not assert convergence to a global optimum of the supervised learning
problem, nor do they imply that every fixed-step or noisy training update is
monotone. Lyapunov stability at strict local minima is a standard property of
smooth gradient systems; it does not identify which basin a given initialization
enters, prove a convergence rate, or analyze the quality of the reached
stationary point.
\end{remark}

\subsection{Complexity}

\begin{proposition}[Complexity per inference step]
\label{prop:complexity}
Let $M=\sum_i |\mathcal{W}(i)|$ be the number of retained motifs and
$S=\sum_i|\mathcal S_g(i)|$ the number of global-support pairs. One inference
step with dense $d\times d$ local/global projections and dense
$\mathbb{R}^{R\times d}$ motif projections costs
\begin{equation}
  O\bigl(|V|Rd^2 + |E|d + MRd + Sd + K|V|d\bigr).
\end{equation}
Under bounded degree $\deg(i)\le\Delta$, one has $M=O(|V|\Delta^2)$; if motif
and global supports are capped by constant budgets per node, the cost is linear
in $|V|$ for fixed $R$, $K$, and support budgets. Dense global support gives
$S=O(|V|^2)$.
\end{proposition}

%----------------------------------------------------------------------
\section{Experimental Protocol}
%----------------------------------------------------------------------

GET should not be evaluated first as a universal graph transformer. The first
test is mechanistic: does FullGET beat a capacity-matched no-motif GET harness
when the label is generated by wedges, triangles, or ternary factors represented
in the motif support? The no-motif harness keeps local pairwise retrieval,
global attention, and HN/FFN memory active while setting $\lambda_3=0$.
Transfer benchmarks are secondary and should narrow the claim when the
mechanism does not transfer.

\paragraph{Protocol principle.}
The core harness is held fixed and the main contrast is always FullGET versus a
capacity-matched no-motif harness ($\lambda_3=0$). Local-only or global-only
variants are ablations, not the primary control. Non-GET structural features,
positional encodings, optimizer family, and wrapper choices are fixed within
each comparison family, while train/validation/test splits are shared across
runs. Later stages count as transfer evidence only if earlier mechanism-level
gains are established.

\paragraph{Stage-gates.}
Stage~0 checks implementation correctness. Stage~1 is decisive for the motif
claim. Stages~2--4 are transfer tests. No Stage~1 gain means no mechanism
claim; no molecular gain means no molecular superiority claim.

\paragraph{Task-conditioned wrappers.}
The energy definition stays fixed. Synthetic and expressivity stages use the
same local/global/HN harness, with global support budgeted if dense attention
would make the diagnostic trivial or expensive. Molecular tasks may add
bond-aware inputs. ET-style transfer may add CLS, LapPE, graph masking, or
stacked blocks only when the comparison family requires them. Wrapper gains
must not be attributed to the motif energy.

\begin{table}[htbp]
\centering
\caption{Experiment families organized by task type, scale, and expected fit
for GET. ``Best fit'' means the experiment can support the main motif-energy
claim; ``diagnostic'' means it is useful but should not carry the headline
claim alone; ``stress/transfer'' means failure narrows the scope rather than
invalidating the mechanism.}
\label{tab:get-task-fit}
\renewcommand{\arraystretch}{1.08}
\footnotesize
\begin{tabularx}{\textwidth}{p{0.24\textwidth} p{0.27\textwidth} p{0.16\textwidth} Y}
\toprule
Family & Scale and signal & GET fit & Claim supported \\
\midrule
Motif-causal synthetic diagnostics
& many small/medium controlled graphs; explicit wedges, triangles, or ternary clauses
& best fit
& isolates whether FullGET beats the no-motif harness because of the order-3 branch \\
Realized expressivity benchmarks
& small/medium structurally hard graphs; indirect local patterns plus symmetry breaking
& diagnostic
& tests whether the learned motif branch realizes useful structural
  discrimination beyond hand-designed labels \\
Bounded-degree molecular graphs
& medium/large collections of small sparse graphs; rings, functional groups, local bond neighborhoods
& best transfer fit
& tests whether local motif gains survive on realistic sparse graph
  prediction tasks \\
Large sparse molecular/peptide graphs
& larger collections or larger per-graph structures; mixed local motifs and longer-range effects
& stress/transfer
& measures scalability and whether motif support remains useful when global
  context also matters \\
TU-style graph classification
& usually small/medium datasets; heterogeneous feature quality and weakly identified motif signal
& secondary transfer
& checks robustness, but gains may come from features, pooling, or wrappers
  rather than motif retrieval \\
Node anomaly graphs
& one or few large graphs; node-level ego-network and homophily/heterophily signals
& ET-adapter stress test
& evaluates graph-domain competitiveness, not the clean higher-order mechanism \\
\bottomrule
\end{tabularx}
\end{table}

\begin{table}[htbp]
\centering
\caption{Stage overview for the compact main-paper protocol.}
\label{tab:stage-wrapper-map}
\renewcommand{\arraystretch}{1.15}
\small
\begin{tabularx}{\textwidth}{l Y Y Y}
\toprule
Stage & Core question & Default setup & Representative tasks \\
\midrule
0 & Is the energy implementation correct? & GET harness on toy graphs; no task-specific wrapper & finite differences, permutation equivariance, monotone Armijo checks \\
1 & Does explicit motif retrieval help when labels are motif-causal? & local/global/HN GET harness with sparse motif support; small/medium controlled graphs & wedge/triangle diagnostics, triangle regression, Max-3-SAT, 3-XORSAT \\
2 & Do Stage~1 gains persist on realized expressivity tests? & GET harness with fixed structural features per method; structurally hard small/medium graphs & BREC by category, CSL, SRG/cycle stress tests \\
3 & Do local gains transfer to bounded-degree molecular graphs? & bond-aware 2D GET; ET-style positional encodings only when required; sparse molecular graph collections & ZINC, ogbg-molhiv, ogbg-molpcba; Peptides as stress test \\
4 & Does an ET-matched wrapper remain competitive on ET-style tasks? & ET-style adapter with CLS/LapPE only when the comparison requires them; heterogeneous benchmark scale & TU-style graph classification, graph anomaly detection, LRGB only with global mixing \\
\bottomrule
\end{tabularx}
\end{table}

\subsection{Stage 1: Mechanism Identification}

Stage~1 is the decisive test of the paper's central claim. We compare
FullGET and a no-motif GET harness on independently sampled graphs drawn from
the same size and degree ranges but differing in wedge or triangle structure,
on triangle-count regression, and on clause-like ternary CSP instances. These
are the best-fit tasks for GET because the support retained by the motif branch
is aligned with the generative factor that determines the label. Difficult
local-structure discrimination tasks such as strongly regular or cycle-parity
style instances are useful stress tests, but they are less direct because
success can depend on propagation depth, symmetry breaking, and structural
encodings. If the motif branch does not beat the capacity-matched no-motif
harness on the primary motif-causal diagnostics, the paper should not claim a
mechanism-level advantage.

Stage~1 should separate oracle support from selected support. Exact
wedge/triangle support tests whether the motif energy can use the right local
factor once supplied; thresholded or budgeted support tests whether the learned
score maps can select useful motifs without giving the model the answer.

\paragraph{Primary motif diagnostics.}
The cleanest Stage~1 tasks are those where the label is generated from the
same local structures represented in the motif support. In wedge
discrimination, graphs are independently sampled under the same size and
degree distribution but are labeled to differ in open-wedge versus
closed-triangle density. In triangle-count regression, the target is a
node-level or graph-level triangle count, with degree-controlled splits so the
model cannot win by learning only a degree surrogate. These tasks directly
test whether the order-3 branch extracts a signal that the no-motif harness lacks under
the same node features, optimizer, inference budget, global support, memory
bank, and readout.

\paragraph{Structural stress tests.}
Strongly regular graph discrimination and cycle-parity detection are useful
stress tests, but they should be reported carefully. Success on a supervised
SRG pair such as Shrikhande versus the $4\times4$ rook graph is evidence that
the trained motif branch realizes a useful structural discriminator; it is not
by itself a formal proof of a global WL-expressivity hierarchy. Cycle-parity
is also a mixed local/global task: triangles are the shortest odd cycles, but
odd-cycle detection in general requires propagation beyond one local motif.
These tasks should therefore be Stage~2-style realized-expressivity
diagnostics or secondary Stage~1 stress tests rather than the only mechanism
evidence.

\subsubsection{Ternary CSP reasoning}
\label{sec:ternary-csp}

Sparse ternary constraint satisfaction is a natural GET task because a
3-clause is already a local order-3 factor. Represent a Max-3-SAT instance as
a signed variable--clause factor graph, keep pairwise variable--clause
incidences in the pairwise branch, and retain exactly three clause-centered
motifs per clause. Let $z_{v_i}=\mathrm{LN}(x_{v_i}^{(T)})$ be the final
variable state and decode a soft assignment by
\begin{equation}
  p_i = \sigma\!\bigl(w_{\mathrm{var}}^\top z_{v_i} + b_{\mathrm{var}}\bigr).
  \label{eq:sat-variable-head}
\end{equation}
For each literal, use $p_i$ or $1-p_i$ according to its sign to obtain a truth
probability $\pi_{ar}$, define clause satisfaction surrogate
$\hat s_a = 1-\prod_{r=1}^{3}(1-\pi_{ar})$, and train with
\begin{equation}
  \mathcal{L}_{\mathrm{Max3SAT}}
  = -\frac{1}{m}\sum_{a=1}^{m}\log(\hat s_a + \varepsilon_{\mathrm{task}}).
  \label{eq:max3sat-loss}
\end{equation}
A 3-XORSAT variant is a useful harder control, but the clean primary
comparison is still no-motif GET versus full GET on the same signed factor
graph. Because the factor graph gives no-motif models the clause--variable
incidences, and the harness retains global attention and HN/FFN memory, the
claimed advantage is not access to clause membership or generic capacity; it is
the explicit ternary interaction among the three variables in a clause.

\paragraph{Rejected Stage~1 candidates.}
Not every task involving triples is a clean test of the current GET core.
Three-coloring is primarily an edge-constraint satisfaction problem; adding a
monochromatic-triangle penalty would test a new hand-designed energy term
rather than the retrieval branch in~\eqref{eq:s3}. Knowledge-graph link
prediction is a relational learning problem with its own negative-sampling and
ranking protocol, so it is better treated as a later relation-aware GET
variant rather than as synthetic mechanism evidence. TSP or shortest-path
tasks with triangle-inequality penalties similarly introduce algorithmic and
metric-consistency objectives outside the present graph-local motif energy.
These tasks may be reasonable extensions, but they should not be used to
support the central Stage~1 claim.

\subsection{Stages 2--4: Transfer}

\paragraph{Stage 2: realized expressivity.}
Use CSL and BREC to test whether gains from Stage~1 survive beyond synthetic
motif-causal labels. The structural-encoding policy for GET should be fixed in
advance; external baselines use their native wrappers, so any positional or
structural features must be reported separately. These datasets are diagnostic
rather than automatically best-fit: they are valuable precisely because they
can expose when local motif retrieval is insufficient without additional
positional or global mixing. BREC should be reported by category rather than
as a single aggregate score. Basic and regular cases can support a
realized-expressivity claim when FullGET improves over the no-motif harness; Extension
and CFI-style failures should be reported as expected limitations of a local
order-3 motif energy rather than hidden inside an average.

\paragraph{Stage 3: bounded-degree molecular graphs.}
Use bounded-degree 2D molecular tasks as the main transfer stage for the core
model. ZINC tests small-graph regression, while ogbg-molhiv and ogbg-molpcba
test whether any local motif benefit survives in more realistic molecular
settings.
Bond-aware inputs are part of the GET wrapper; stronger positional encodings
are used only when the comparison family requires them, and they should be
reported separately from motif retrieval. If the motif branch does not help on
these tasks, the empirical claim should remain limited to mechanism and
expressivity benchmarks. This is the most plausible real-data regime for GET
because sparse molecular graphs contain many bounded local substructures, but
the comparison must still separate gains from bond features, structural
encodings, and motif retrieval itself. Peptides and other LRGB tasks are
useful stress tests, but they should be marked as long-range benchmarks; poor
performance there does not by itself refute the local motif mechanism, and
strong performance should not be claimed without controlling for global mixing
and positional encodings. In practice, this stage should be the default
optimization target for systems work: molecule batches are often small enough
for efficient sparse local energy kernels while still giving the clearest
transfer signal for local higher-order terms.

\paragraph{Stage 4: ET-style transfer.}
ET-style anomaly detection and graph-classification benchmarks are reported as
transfer tests rather than as the primary proof of the higher-order mechanism.
When these comparisons are run, GET should inherit the same ET-matched wrapper
components that the benchmark rewards---for example CLS readout, Laplacian
positional encodings, or graph-aware masking---instead of attributing wrapper
gains to the motif term itself. TU-style graph classification and node anomaly
benchmarks are therefore useful for external validity and engineering
readiness, but they are not the cleanest evidence for the proposed order-3
energy because dataset size, feature quality, class imbalance, and wrapper
choices can dominate the result.

\paragraph{Closest external comparators.}
The current repository implements pairwise controls such as GIN, GCN, GAT,
GraphTransformer, and BWGNN, plus ET-style adapters. Higher-order baselines
such as GSN/ESAN-style substructure methods, k-GNN/K-hop variants, simplicial
message passing, and Triplet Graph Transformer are required external
comparators for a submission-level motif claim. If they are not implemented, the
evidence should be described only as a controlled no-motif GET ablation, not as
competitiveness against higher-order graph models. Graphormer, GPS, GRIT, and
Exphormer remain transfer comparators for graph-transformer harnesses.

\paragraph{Baselines and ablations.}
The compact ablation suite should include local-only GET, local+global GET,
local+global+HN no-motif GET, Motif-only GET, FullGET, open-wedge versus
open+closed support, rank-$1$ versus rank-$R$ motif factorization, fixed-step
versus Armijo inference, global-support budget, and HN/FFN on/off. The primary
comparison is always FullGET versus a capacity-matched no-motif GET harness;
parameter counts should be matched to within a small tolerance so that raw
width does not masquerade as a motif gain.

\paragraph{ET comparison discipline.}
ET-style models are pairwise in their attention energy, but the manuscript
should not claim that ET ``cannot solve'' a synthetic task in an absolute
sense. Dense attention, CLS tokens, positional encodings, node identifiers, or
hand-crafted structural features can let a pairwise transformer fit many
finite synthetic datasets. The fair claim is narrower: under matched inputs
and without explicit triangle, clause, or motif features, an ET-style no-motif
GET control lacks the local order-3 retrieval term that FullGET uses.
The evidence for GET is therefore a controlled performance gap against that
matched no-motif control, not a blanket impossibility statement about ET or
graph transformers.

\paragraph{Model selection and reporting.}
Stage~0 should verify finite-difference gradients, permutation equivariance,
and monotone decrease under accepted Armijo steps on random toy graphs.
Hyperparameter tuning should be stage-gated and non-leaky: shortlist
configurations on the primary Stage~1 diagnostics, carry them to later stages
unchanged, and report mean and standard deviation across multiple seeds
together with paired comparisons against the strongest matched baseline. Each
Stage~1 task should include energy-decay curves, performance versus inference
steps, and ablations over $\lambda_3$, motif rank $R$, motif budget,
global-support budget, and motif-type vocabulary. Runnable baselines in this
repository are pairwise controls such as GIN, GCN, GAT, and GraphTransformer,
plus BWGNN and ET-style adapters when the benchmark is ET-like. Higher-order
families such as GSN-style models, k-GNN/K-hop variants, simplicial message
passing, Graphormer, GPS, GRIT, and Exphormer remain future comparison targets
unless they are added explicitly. This makes it possible to separate
improvements from higher-order energy terms, from global/HN harness capacity,
from static motif features, from generic structural encodings, and from
task-specific wrappers.

%----------------------------------------------------------------------
\section{Discussion}
%----------------------------------------------------------------------

\paragraph{What is and is not claimed.}
The entropy-regularized lifting is a standard log-sum-exp identity. It is
included to define the routing variables and expose the full-gradient update,
not as a new variational inference method. GET is not BP, GBP, Kikuchi
inference, VMP, the first energy-based GNN, or the first Hopfield model on
graphs. The claim is narrower: explicit sparse motif retrieval is placed inside
one local/global/HN Hopfield/ET-style graph energy and tested against the
matched no-motif energy harness.

\paragraph{Limitations.}
The monotone-descent result concerns accepted inference steps with fixed
parameters. It does not imply faster supervised training, global convergence,
or better generalization. The base motif branch is symmetric in the two neighbor
roles and may underspecify directed or role-asymmetric interactions. Strong
structural encoders, persistent edge states, geometric supervision, or triplet
transformer architectures may capture similar signals with better trade-offs on
some tasks \citep{ying2021graphormer,rampasek2022gps,ma2023grit,
shirzad2023exphormer,canturk2024gpse,airale2025spse,diao2023relational,
hussain2024triplet}.

\paragraph{Task-conditioned variants.}
The core GET mechanism should be presented as a family of task-conditioned
instantiations rather than as one frozen checkpoint. Synthetic tasks use the
same local/global/HN energy harness with controlled support budgets;
realized-expressivity tasks keep that harness but fix the structural-encoding
policy explicitly; molecular tasks add bond-aware inputs when available; and
ET-style transfer uses an ET-faithful adapter when CLS readout, LapPE,
graph-aware masking, or stacked blocks are part of the comparison family. This
separation helps keep wrapper choices from being conflated with the motif
energy itself.

\paragraph{Extensions.}
The present model covers local retrieval, motif retrieval, global attention,
and HN/FFN memory. Further extensions include dynamic edge states,
task-conditioned readout attractors, typed role-asymmetric motifs, structural
memory latents, and joint multi-layer HAM-style energies. Those additions
should be named GET+X variants until implemented and ablated.

%----------------------------------------------------------------------
\section{Conclusion}
%----------------------------------------------------------------------

GET is a scalar-energy graph-transformer model with graph-local pairwise
retrieval, explicit sparse motif retrieval, ET-style global attention, and
HN/FFN memory. The log-sum-exp terms give entropy-regularized routing weights,
while the full gradient gives ET-style anchor and reverse-incidence messages.
The paper's test is narrow: establish whether the motif energy improves over a
matched no-motif energy harness on motif-causal tasks, then report how far that
mechanism transfers.

\appendix
%----------------------------------------------------------------------
\section{LayerNorm Jacobian}
\label{app:layernorm}
%----------------------------------------------------------------------

For $d\ge2$, standard $\varepsilon$-stabilized LayerNorm is
$\mathrm{LN}(x) = a_{\mathrm{LN}} \odot \hat{x} + b_{\mathrm{LN}}$, where
$\hat{x} = (x - \mu\mathbf{1})/\sigma$,
$\mu = \frac{1}{d}\mathbf{1}^\top x$, and
$\sigma = \sqrt{\frac{1}{d}\|x - \mu\mathbf{1}\|^2 + \varepsilon}$.
For $\varepsilon > 0$, $\sigma > 0$ everywhere, so $\mathrm{LN}$ is
$C^\infty$. The Jacobian is \citep{ba2016layer}:
\begin{equation}
  J_{\mathrm{LN}}(x)
  = \frac{1}{\sigma}\mathrm{Diag}(a_{\mathrm{LN}})
    \Bigl(I - \tfrac{1}{d}\mathbf{1}\mathbf{1}^\top
           - \tfrac{1}{d}\hat{x}\hat{x}^\top\Bigr).
\end{equation}

%----------------------------------------------------------------------
\section{Lyapunov Stability Derivation}
\label{app:lyapunov}
%----------------------------------------------------------------------

We give the detailed derivation of Proposition~\ref{prop:lyapunov}. Let
$X^*$ be a strict local minimizer of $E_{\mathrm{GET}}$ and define the shifted
energy
\begin{equation}
  V(X) = E_{\mathrm{GET}}(X) - E_{\mathrm{GET}}(X^*).
\end{equation}
Then $V(X^*)=0$. Because $X^*$ is a strict local minimizer and
$E_{\mathrm{GET}}$ is continuous, there exists $r_0>0$ such that
\begin{equation}
  0 < \|X-X^*\|_F \le r_0
  \quad\Longrightarrow\quad
  V(X) > 0.
  \label{eq:lyapunov-local-positive}
\end{equation}
Thus $V$ is locally positive definite around $X^*$.

To prove Lyapunov stability, fix any $\varepsilon\in(0,r_0)$. Since the
sphere
\begin{equation}
  S_\varepsilon = \{X : \|X-X^*\|_F = \varepsilon\}
\end{equation}
is compact and $V$ is continuous and strictly positive on it by
\eqref{eq:lyapunov-local-positive}, the minimum
\begin{equation}
  \alpha_\varepsilon
  = \min_{X\in S_\varepsilon} V(X)
  > 0
  \label{eq:lyapunov-alpha}
\end{equation}
is attained and strictly positive. By continuity of $V$ at $X^*$, there
exists $\delta\in(0,\varepsilon)$ such that
\begin{equation}
  \|X-X^*\|_F < \delta
  \quad\Longrightarrow\quad
  V(X) < \alpha_\varepsilon.
  \label{eq:lyapunov-delta}
\end{equation}
Consequently, the local sublevel set
\begin{equation}
  \mathcal{S}_{\alpha_\varepsilon}
  = \{X : V(X) < \alpha_\varepsilon\}
  \label{eq:lyapunov-sublevel}
\end{equation}
is contained in the open ball $B_\varepsilon(X^*)$.

\paragraph{Continuous-time gradient flow.}
Consider the inference ODE
\begin{equation}
  \dot X = -\nabla_X E_{\mathrm{GET}}(X).
\end{equation}
Because $E_{\mathrm{GET}}$ is $C^1$, the chain rule gives
\begin{equation}
  \dot V(X)
  = \bigl\langle \nabla_X V(X), \dot X \bigr\rangle
  = -\|\nabla_X E_{\mathrm{GET}}(X)\|_F^2
  \le 0.
  \label{eq:lyapunov-flow-decay}
\end{equation}
Hence $V$ is nonincreasing along any gradient-flow trajectory.

\paragraph{Discrete-time Armijo descent.}
For the accepted Armijo step
\begin{equation}
  X^{t+1} = X^t - \eta_t\nabla_X E_{\mathrm{GET}}(X^t),
\end{equation}
Theorem~\ref{thm:monotone} gives
\begin{equation}
  V(X^{t+1})
  \le V(X^t) - c\eta_t\|\nabla_X E_{\mathrm{GET}}(X^t)\|_F^2
  \le V(X^t).
  \label{eq:lyapunov-armijo-decay}
\end{equation}

\paragraph{Invariant-neighborhood argument.}
If $\|X(0)-X^*\|_F<\delta$, then $V(X(0))<\alpha_\varepsilon$. Since $V$ is
nonincreasing along gradient flow, the trajectory cannot reach
$S_\varepsilon$, because every point on that sphere has
$V\ge\alpha_\varepsilon$. Thus $\|X(t)-X^*\|_F<\varepsilon$ for all $t\ge0$.
The same contradiction argument applies to Armijo iterates using
\eqref{eq:lyapunov-armijo-decay}. Hence $X^*$ is Lyapunov stable for both
dynamics.

\begin{corollary}[Local asymptotic stability of the gradient flow]
\label{cor:lyapunov-asymptotic}
Assume the conditions of Proposition~\ref{prop:lyapunov}. Suppose in addition
that there exists $r>0$ such that $X^*$ is the only stationary point of
$E_{\mathrm{GET}}$ in the ball $B_r(X^*)$. Then $X^*$ is locally
asymptotically stable for the continuous-time gradient flow
$\dot X = -\nabla_X E_{\mathrm{GET}}(X)$.
\end{corollary}

\begin{proof}
By Proposition~\ref{prop:lyapunov}, $X^*$ is Lyapunov stable. Choose
$\varepsilon>0$ small enough that $B_\varepsilon(X^*)\subset B_r(X^*)$.
Inside this ball, $\dot V=-\|\nabla_X E_{\mathrm{GET}}(X)\|_F^2$. The largest
invariant subset of $\{\dot V=0\}$ is therefore the set of stationary points,
which by assumption is $\{X^*\}$. The local invariance principle gives
$X(t)\to X^*$ for all trajectories initialized sufficiently close to $X^*$.
\end{proof}

\begin{corollary}[Positive-definite Hessian as a sufficient condition]
\label{cor:lyapunov-hessian}
Assume the standing smoothness conditions of Proposition~\ref{prop:lipschitz},
and suppose additionally that the Euclidean Hessian of $E_{\mathrm{GET}}$ with
respect to $\operatorname{vec}(X)$ is positive definite at $X^*$:
\begin{equation}
  \nabla^2_{\operatorname{vec}(X)} E_{\mathrm{GET}}(X^*) \succ 0.
\end{equation}
Then $X^*$ is a strict local minimizer and an isolated stationary point.
Consequently, $X^*$ is locally asymptotically stable for the continuous-time
gradient flow.
\end{corollary}

\begin{proof}
Positive definiteness of the Hessian gives the standard second-order
sufficient condition for a strict local minimum and implies that stationary
points cannot accumulate at $X^*$. Proposition~\ref{prop:lyapunov} gives
Lyapunov stability, and Corollary~\ref{cor:lyapunov-asymptotic} gives local
asymptotic stability for the gradient flow.
\end{proof}

\begin{remark}
The asymptotic corollaries are stated only for the continuous-time gradient
flow. A corresponding local asymptotic-stability statement for accepted
Armijo iterates would require additional discrete-time assumptions, such as a
uniform lower bound on accepted step sizes or a separate convergence theorem.
\end{remark}

%----------------------------------------------------------------------
\section{Gradient Derivations}
\label{app:gradients}
%----------------------------------------------------------------------

The following formulas differentiate the log-sum-exp terms directly. Equivalently,
by Danskin's theorem applied to Proposition~\ref{prop:routing}, the derivative
with respect to each score is the optimal routing weight; no additional
softmax-Jacobian terms appear in the score gradient.

Define softmax weights:
\begin{equation}
  \alpha_{ij}
  =
  \begin{cases}
    \dfrac{\exp(\beta_2 s^{(2)}_{ij})}
          {\sum_{\ell \in \mathcal{N}(i)} \exp(\beta_2 s^{(2)}_{i\ell})},
      & \mathcal{N}(i)\neq\varnothing,\ j \in \mathcal{N}(i), \\[1em]
    0, & \text{otherwise,}
  \end{cases}
\end{equation}
\begin{equation}
  \gamma_{i,jk}
  =
  \begin{cases}
    \dfrac{\exp(\beta_3 s^{(3)}_{i,jk})}
          {\sum_{(u,v) \in \mathcal{W}(i)} \exp(\beta_3 s^{(3)}_{i,uv})},
      & \mathcal{W}(i)\neq\varnothing,\ (j,k)\in\mathcal{W}(i), \\[1em]
    0, & \text{otherwise,}
  \end{cases}
\end{equation}
\begin{equation}
  \rho_{ia}
  = \frac{\exp(\beta_m s^{(m)}_{i,a})}
         {\sum_{c=1}^{K} \exp(\beta_m s^{(m)}_{i,c})}.
\end{equation}
\begin{equation}
  \pi_{ij}
  =
  \begin{cases}
    \dfrac{\exp(\beta_g s^{(g)}_{ij})}
          {\sum_{\ell \in \mathcal{S}_g(i)} \exp(\beta_g s^{(g)}_{i\ell})},
      & \mathcal{S}_g(i)\neq\varnothing,\ j \in \mathcal{S}_g(i), \\[1em]
    0, & \text{otherwise.}
  \end{cases}
\end{equation}

The full gradient with respect to the raw state $x_i$ is
\begin{equation}
  \nabla_{x_i} E_{\mathrm{GET}}(X)
  = x_i + {J_{\mathrm{LN}}(x_i)}^\top
    \bigl(\nabla_{g_i} E_{\mathrm{att}} + \nabla_{g_i} E_{\mathrm{HN}}\bigr).
\end{equation}

\paragraph{Attention branch.}
Define the motif composite key
\begin{equation}
  \phi_{ijk}^{(r)}
  = k_3^{(r)}(g_j) \odot k_3^{(r)}(g_k) + t_{\tau_{ijk}}^{(r)}.
\end{equation}
Suppressing the rank superscript in the Jacobian notation, the gradient is
\begin{equation}
\begin{aligned}
  \nabla_{g_i} E_{\mathrm{att}}
  ={} &
  -\frac{\lambda_2}{\sqrt{d}}
  \Bigl[
    {J_{q_2}(g_i)}^\top \sum_{j \in \mathcal{N}(i)} \alpha_{ij}\,k_2(g_j)
    + {J_{k_2}(g_i)}^\top \sum_{u:\,i \in \mathcal{N}(u)} \alpha_{ui}\,q_2(g_u)
  \Bigr] \\
  &-\frac{\lambda_3}{\sqrt{Rd}}
  \Bigl[
    \sum_{r=1}^R {J_{q_3^{(r)}}(g_i)}^\top
      \sum_{(j,k)\in\mathcal{W}(i)} \gamma_{i,jk}\,\phi_{ijk}^{(r)}
    + \sum_{r=1}^R {J_{k_3^{(r)}}(g_i)}^\top \Psi_i^{(r)}
  \Bigr]\\
  &-\frac{\lambda_g}{\sqrt{d}}
  \Bigl[
    {J_{q_g}(g_i)}^\top \sum_{j \in \mathcal{S}_g(i)} \pi_{ij}\,k_g(g_j)
    + {J_{k_g}(g_i)}^\top \sum_{u:\,i \in \mathcal{S}_g(u)} \pi_{ui}\,q_g(g_u)
  \Bigr],
\end{aligned}
\end{equation}
where $\Psi_i^{(r)}$ collects motifs in which node $i$ appears as a key:
\begin{equation}
  \Psi_i^{(r)}
  = \sum_{u \in V}\; \sum_{\substack{k:\,(i,k)\in\mathcal{W}(u)}}
      \gamma_{u,ik}\bigl(q_3^{(r)}(g_u) \odot k_3^{(r)}(g_k)\bigr)
  +
    \sum_{u \in V}\; \sum_{\substack{j:\,(j,i)\in\mathcal{W}(u)}}
      \gamma_{u,ji}\bigl(q_3^{(r)}(g_u) \odot k_3^{(r)}(g_j)\bigr).
\end{equation}
The two sums handle the cases where $i$ appears as the first or second element
of the stored representative of an unordered retained pair; since
$\mathcal W(u)$ stores each pair once, this does not double-count a motif.

\paragraph{Memory branch.}
\begin{equation}
  \nabla_{g_i} E_{\mathrm{HN}}
  = -\frac{\lambda_m}{\sqrt{d}}\,
    {J_{q_m}(g_i)}^\top \sum_{a=1}^K \rho_{ia}\,k_m(b_a).
\end{equation}

If the memory vectors are trained, their gradients are
\begin{equation}
  \nabla_{b_a} E_{\mathrm{HN}}
  = -\frac{\lambda_m}{\sqrt{d}}\,
    {J_{k_m}(b_a)}^\top \sum_{i \in V} \rho_{ia}\,q_m(g_i).
\end{equation}

\bibliographystyle{plainnat}
\bibliography{ref}

\end{document}
